\section{Extension to H-index}
\label{sec:hindex}

In this section, we extend the sketch propagation framework to approximately process \ced queries based on the h-index.
To compute the h-index of each node, we should compute the degrees of its neighbors.
However, KMV sketches can only estimate the neighborhood sizes of a node but fail to provide the degree estimations of its neighbors.
Therefore, we propose a novel pivot-based algorithm to obtain the set of hubs in terms of h-index without explicitly calculating the h-index of every node in $H$.
Before delving into the algorithm, we first review the definition of \emph{h-index} \cite{LuZZS16}.
\begin{definition}[h-index]
  Given a node $u \in H$, the h-index of $u$ is the largest integer $h(u)$ such that $u$ has no fewer than $h(u)$ neighbors of degrees at least $h(u)$.
\end{definition}

\subsection{Pivot-based Algorithm for \ced}
\label{sec:hindex-global}

The basic idea of the pivot-based algorithm is to choose a random pivot node $u \in V$ and iteratively partition the nodes in $H$ into two disjoint sets, $Q^+_u$ and $Q^-_u$, based on their h-index values in comparison with $h(u)$, i.e., $Q^+_u = \{v \in V \,|\, h(v) \geq h(u)\}$ and $Q^-_u = \{v \in V \,|\, h(v) < h(u)\}$.
As such, we can decide that $Q^+_u \subseteq S^+_0(v_\lambda)$ if $u \in S^+_0(v_\lambda)$ since all nodes in $Q^+_u$ have h-index values no less than $h(u)$ and thus are ranked higher than $u$, or $Q^-_u \cap S^+_0(v_\lambda) = \varnothing$ if $u \in S^-_0(v_\lambda)$ similarly.
This partitioning strategy allows for efficient bisection when identifying hubs, i.e., $S^+_0(v^\lambda)$.
If $u \in S^+_0(v_\lambda)$, we can add all nodes in $Q^+_u$ to the result set and exclude them from consideration subsequently.
In contrast, if $u \in S^-_0(v^\lambda)$, the nodes in $Q^+_u$ may contain those in $S^+_0(v^\lambda)$ and $S^-_0(v^\lambda)$, which require further partitioning in subsequent iterations. Meanwhile, all nodes in $Q^-_u$ must not be in $S^+_0(v^\lambda)$ and are excluded from consideration.
The iterative process above proceeds until all nodes have been decided.
Then, all nodes in $S^+_0(v^\lambda)$ are added to the result.

\begin{algorithm}[t]
\caption{Pivot-based Algorithm}
\label{alg:hindex}
\LinesNumbered
\KwIn{KG $\gset$, meta-path $\mset$, quantile parameter $\lambda$, number of propagations $\theta$, KMV sketch size $k$}
\KwOut{The set of hubs $S$}
    \ForAll{$u \in \vset^*$}{ \label{algline:pivotstagei-begin}
        \For{$t=1$ {\rm \textbf{to}} $\theta$}{
            $\kset[u][t] \gets \varnothing$\;
            \textbf{if} $u \in \vset_0$ \textbf{do} add $r(u) = \mathtt{Rand}(0, 1)$ to $\kset[u][t]$\;
        }
    }
    $\widetilde{N} \gets$ \texttt{Propagate}$(\kset)$\; \label{algline:pivotstagei-end}
    $Q \gets V$ and $S \gets \varnothing$\; 
    \While{$|Q|>0$}{ \label{algline:iterQ-begin}
        Sample a random node from $Q$ as the pivot node $u$\; \label{algline:hest-begin}
        Sort nodes in $\nset(u)$ in descending order by $\widetilde{N}$\;
        Initialize $\widetilde{h}(u) \gets 0$\;
        \ForAll{$v \in \nset(u)$ and $\widetilde{N}(v) \geq \widetilde{h}(u)$}{
            $\widetilde{h}(u) \gets \widetilde{h}(u) + 1$\; \label{algline:hest-end}
        }
        \ForAll{$v \in \vset^*$}{ \label{algline:initK2-begin}
            \For{$t=1$ \KwTo $\theta$}{
                $\kset[v][t] \gets \varnothing$\;
                \If{$v \in \vset_0$ and $\widetilde{N}(v) \geq \widetilde{h}(u)$}{
                    Add $r(v) = \mathtt{Rand}(0, 1)$ to $\kset[v][t]$\; \label{algline:initK2-end}
                }
            }
        }
        $\widetilde{h} \gets$ \texttt{Propagate}$(\kset)$\; \label{algline:prop2}
        Initialize $Q^+_u \gets \varnothing$ and $Q^-_u \gets \varnothing$\; \label{algline:partitionQ-begin}
        \ForAll{$v\in Q$}{
            \textbf{if} $\widetilde{h}(v) \geq \widetilde{h}(u)$ \textbf{do} add $v$ to $Q^+_u$\; 
            \textbf{else} add $v$ to $Q^-_u$\; \label{algline:partitionQ-end}
        }
        \uIf{$|Q^+_u| + |S| \leq \lambda|V|$}{ \label{algline:updateS-begin}
            $Q \gets Q^-_u$ and $S \gets S \cup Q^+_u$\;
            \textbf{if} $|S| \leq \lambda|V|$ \textbf{then} $S \gets S \cup \{u\}$\;
        }
        \textbf{else} $Q\gets Q^+_u$\; \label{algline:updateS-end}
    } \label{algline:iterQ-end}
    \Return{$S$}\; \label{algline:pivotreturn}
\end{algorithm}

\vspace{1mm}
\myp{Pivot-based Partitioning}
The key challenge lies in partitioning the nodes in $H$ without explicitly computing the h-index of each node.
In fact, we can compare the h-indexes of two nodes using only one of them.
By the definition of h-index, for any $u, v \in V$, we have $h(v) \geq h(u)$ iff $v$ has at least $h(u)$ neighbors in $H$ with degrees larger than or equal to $h(u)$.
Based on the above observation, we introduce a two-stage pivotal partitioning method for our pivot-based algorithm.

\noindent\underline{\textit{Stage (I): Estimate $h(u)$.}}
We first employ \emph{sketch propagation} on the matching graph to estimate the degrees of all nodes in $H$.
Subsequently, at each iteration, we scan the set of neighbors $\nset(u)$ of a randomly chosen node $u$ (obtained by a single DFS on the matching graph starting from $u$).
Based on the degree estimations, we compute an approximate h-index $\widetilde{h}(u)$ of node $u$.

\noindent\underline{\textit{Stage (II): Estimate $Q^+_u$ and $Q^-_u$.}}
Once we have the approximate h-index $\widetilde{h}(u)$, we proceed to eliminate the nodes in $H$ with degrees smaller than $\widetilde{h}(u)$.
Next, we perform \emph{sketch propagation} on the subgraph of the matching graph that only includes the remaining nodes.
The resulting KMV sketches estimate the degrees of nodes in the subgraph of $H$ induced by the set of nodes with degrees greater than or equal to $\widetilde{h}(u)$ in the original graph $H$, denoted as $\widetilde{H}(u)$.
Finally, we obtain $Q^+_u$ as the set of nodes with estimated degrees greater than $\widetilde{h}(u)$ in the induced subgraph of $\widetilde{H}(u)$, while the remaining nodes are added to $Q^-_u$.

Algorithm~\ref{alg:hindex} depicts the procedure of our pivot-based algorithm.
It also receives a KG $\gset$, a meta-path $\mset$, the quantile parameter $\lambda$, the number of propagations $\theta$, and the KMV sketch size $k$ as input, and returns a set of nodes $S$ to answer the (approximate) \ced query based on h-index.
It first calls the same \texttt{Propagate} function as Algorithm~\ref{alg:propagation} over the array $\kset$ of KMV sketches to estimate the degree of each node in $H$ (Lines~\ref{algline:pivotstagei-begin}--\ref{algline:pivotstagei-end}).
Note that the estimated degrees will be reused across all iterations for partitioning.
The iterative process proceeds until $Q$, which contains the nodes to be partitioned in each iteration, is empty (Lines~\ref{algline:iterQ-begin}--\ref{algline:iterQ-end}).
In each iteration, it first randomly chooses a pivot node $u$ from $Q$ and estimates its h-index $\widetilde{h}(u)$ (Lines~\ref{algline:hest-begin}-\ref{algline:hest-end}).
Then, it re-initializes arrays $\kset$ of KMV sketches for the nodes with degrees greater than $\widetilde{h}(u)$ (Lines~\ref{algline:initK2-begin}--\ref{algline:initK2-end}).
After that, the \texttt{Propagate} function is performed over $\kset$ to estimate the number of neighbors with degrees greater than $\widetilde{h}(u)$ for each node in $Q$ (Line~\ref{algline:prop2}).
Accordingly, it divides $Q$ into $Q^+_u$ and $Q^-_u$ according to $\widetilde{h}$ (Lines~\ref{algline:partitionQ-begin}--\ref{algline:partitionQ-end}) and updates the result set $S$ and the candidate node set $Q$ based on $Q^+_u$ and $Q^-_u$ (Lines~\ref{algline:updateS-begin}--\ref{algline:updateS-end}).
Finally, when $Q$ is empty, the result set $S$ is returned (Line~\ref{algline:pivotreturn}). The ties are broken arbitrarily.

\myp{Theoretical Analysis}
Algorithm~\ref{alg:hindex} returns the correct answer for any h-index-based \ced query w.h.p.
We first give the following lemma, which states that Algorithm~\ref{alg:hindex} provides a concentrated estimation of $h(u)$ for the randomly chosen pivot node $u$ in each iteration.
\begin{lemma}\label{lemma:hq}
For any $\delta, p\in(0, 1)$, if $\theta=\Theta(\frac{\Lambda}{\delta k}\log{\frac{|V|}{p}})$, we have $(1-\delta) \cdot h(u) \leq \widetilde{h}(u) \leq (1+\delta) \cdot h(u)$ with probability at least $1 - p$ for the pivot node $u$ in each iteration.
\end{lemma}

\techreport{
\begin{proof}
Let us order the nodes in $\nset(u)$ according to their degrees in $H$ descendingly and denote $v_i$ as the $i$-th neighbor of $u$.
Then, we have $N(v_i) \geq h(u)$ for $0 \leq i \leq h(u) - 1$ and $N(v_i) \leq h(u)$ for $h(u) \leq i \leq N(u) - 1$.
By combining Lemma~\ref{theorem:dconcentration} with the union bound over $V$, when $\theta = \Theta(\frac{\Lambda}{\delta k}\log{\frac{|V|}{p}})$ in the first stage, we have $\widetilde{N}(v_i) \geq (1-\delta)N(v_i) \geq (1-\delta)h(u)$ for $0 \leq i \leq h(u)-1$ with probability at least $1 - p$.
Thus, node $u$ has $h(u) \geq (1-\delta) h(u)$ neighbors with estimated degrees of at least $(1-\delta) h(u)$.
Thus, $\widetilde{h}(u) \geq (1-\delta) h(u)$ with probability at least $1 - p$.
Similarly, we have $\widetilde{N}(v_i) \leq (1+\delta)N(v_i) \leq (1+\delta) h(u)$ for $i \geq h(u)$ with probability at least $1 - p$.
Thus, node $u$ has no more than $h(u)$ neighbors with estimated degrees greater than $(1+\delta) h(u)$, and $\widetilde{h}(u) \leq (1+\delta) h(u)$ with probability at least $1 - p$.
\end{proof}
}

The following lemma indicates that the pivot-based partitioning method divides candidate nodes into $\epsilon$-separable sets. 
\begin{lemma}\label{lemma:concenpar}
  For any $\delta^\prime \in (0, 1)$ and pivot node $u$, by setting $\theta = \Theta(\frac{\Lambda}{\delta^\prime k} \log{\frac{|V|}{p}})$, we have $S^-_{\delta^\prime}(u) \subseteq Q^-_u$ and $S^+_{\delta^\prime}(u) \subseteq Q^+_u$ with probability at least $1 - p$.
\end{lemma}

\techreport{
\begin{proof}
  For any node $v \in S^-_{\delta^\prime}(u)$, let $\nset_u(v)$ denote the neighbors of $v$ in $\widetilde{H}(u)$ and $N_u(v)=|\nset_u(v)|$.
  Since $h(v) < (1-\delta^\prime) h(u)$, $v$ has at most $(1-\delta^\prime) h(u)$ neighbors with degrees larger than or equal to $(1-\delta^\prime) h(u)$.
  For each neighbor $v^\prime$ of $v$ with $N(v^\prime) < (1-\delta^\prime) h(u)$, when $\theta = \Theta(\frac{\Lambda}{\delta k} \log{\frac{|V|}{p}})$, $\widetilde{N}(v^\prime) < (1+\delta)N(v^\prime) < (1+\delta)(1-\delta^\prime) h(u) < \frac{(1+\delta)(1-\delta^\prime)}{1-\delta} \widetilde{h}(u)$.
  By setting $\delta < \frac{\delta^\prime}{2-\delta^\prime}$, we have $\widetilde{N}(v^\prime) < \widetilde{h}(u)$, i.e., $v^\prime$ is filtered out in \textit{Stage (I)}, with probability at least $1 - p$.
  Thus, we have $N_u(v) < (1-\delta^\prime) h(u)$ with probability at least $1 - p$.
  Let $\widetilde{N}_u(v)$ denote the estimation of $N_u(v)$ in \textit{Stage (II)} of pivot-based partitioning.
  We have $\widetilde{N}_u(v) \leq (1+\delta) N_u(u)$ with probability at least $1 - p$, and thus $\widetilde{N}_u(v) \leq (1+\delta) N_u(v) < (1+\delta)(1-\delta^\prime) h(u) \leq \frac{(1+\delta)(1-\delta^\prime)}{1-\delta} \widetilde{h}(u) \leq \widetilde{h}(u)$, i.e., $v$ is added to $Q^-_u$ with probability at least $1 - p$.
  Thus, by setting $\theta = \Theta(\frac{\Lambda}{\frac{\delta^\prime}{2-\delta^\prime} k} \log{\frac{|V|}{p}}) = \Theta(\frac{\Lambda}{\delta^\prime k} \log{\frac{|V|}{p}})$, we have $S^-_{\delta^\prime}(u) \subseteq Q^-_u$.
  Similarly, we also get $S^+_{\delta^\prime}(u) \subseteq Q^+_u$ with probability at least $1 - p$ when $\theta = \Theta(\frac{\Lambda}{\delta^\prime k} \log{\frac{|V|}{p}})$.
\end{proof}
}

Subsequently, we can formally analyze the correctness of the pivot-based algorithm for processing $\epsilon$-approximate \ced queries based on the h-index in the following theorem.
\begin{theorem}\label{theorem:hindex}
  Let $\theta = \Theta(\frac{\Lambda}{\epsilon k} \log{\frac{|V|}{p}})$. Algorithm~\ref{alg:hindex} returns the answer to an $\epsilon$-approximate \ced query based on the h-index correctly with probability at least $1 - p$.
\end{theorem}

\techreport{
\begin{proof}
  We prove the theorem by examining three different cases for the ranges of the h-index $h(u)$ for the pivot node $u$.
  For each case, we show the correctness of the first iteration in Algorithm~\ref{alg:hindex}, while subsequent iterations are proven by induction.
  By selecting $\delta^\prime < \frac{\epsilon}{2}$ as Lemma~\ref{lemma:concenpar}, we have:
  \begin{itemize}[leftmargin=*]
    \item \textbf{Case 1 ($h(u) > (1+\frac{\epsilon}{2}) h(v^\lambda)$):} Each $v$ with $h(v)\leq h(v^\lambda)$ will be added to $Q^-_u$.
    Thus, $|Q^-_u| \geq (1-\lambda) |V|$.
    Therefore, $Q^-_u$ will be used as the candidate set $Q$ in the next iteration, and $Q^+$, which does not contain any node in $S^-_\epsilon(v^\lambda)$, will be added to $S$.
    \item \textbf{Case 2 ($h(u) < (1-\frac{\epsilon}{2}) h(v^\lambda)$):} Each $v$ with $h(v)\geq h(v^\lambda)$ will be added to $Q^+_u$.
    Thus, $|Q^+_u| \geq \lambda |V|$.
    Therefore, $Q^+$, which contains all nodes in $S^+_\epsilon(v^\lambda)$, will be used as the candidate set $Q$ in the next iteration.
    \item \textbf{Case 3 ($(1-\frac{\epsilon}{2}) h(v^\lambda) < h(u) < (1+\frac{\epsilon}{2}) h(v_\lambda)$):} Any $v$ with $h(v) \geq (1+\epsilon) h(v^\lambda)$ will be added to $Q^+_u$.
    And any $v^\prime$ with $h(v^\prime) \leq (1-\epsilon) h(v^\lambda)$ will be added to $Q^-_u$.
    Thus, $Q^+_u$, which contains all nodes in $S^+_\epsilon(v^\lambda)$, will be added to $S$ or used as the candidate set $Q$ in the next iteration; whereas $Q^-_u$, which contains all nodes in $S^-_\epsilon(v^\lambda)$, will be eliminated or used as the candidate set $Q$ for the next iteration.
  \end{itemize}
  Considering all of the above cases, it follows that, for any randomly chosen pivot $u$, none of the nodes in $S^-_\epsilon(v^\lambda)$ is added to $S$. In contrast, all nodes in $S^+_\epsilon(v^\lambda)$ have a probability of at least $1 - p$ to be either added to $S$ or retained for the subsequent iteration.
  Thus, we prove the theorem by applying the union bound over all iterations.
\end{proof}
}

Finally, the amortized time complexity of Algorithm~\ref{alg:hindex} is $O(\frac{\Lambda |\eset^*|}{\epsilon} \cdot \log^2{\frac{|V|}{p}})$ because it invokes \emph{sketch propagation} $O(\log{|V|})$ times in amortization. Its space complexity is $O(k\theta|\vset^*|+|\eset^*|)$, which is equal to that of Algorithm~\ref{alg:propagation}.

\subsection{Early Termination for Personalized \ced}
\label{sec:hindex-personal}

Given a query node $q$, the personalized \ced query based on the h-index can also be answered by performing the pivot-based partitioning method with $q$ as the pivot node to obtain $Q^+_u$ and returning \texttt{FALSE} if $|Q^+_u| > \lambda|V|$ or \texttt{TRUE} otherwise.
Following Lemma~\ref{lemma:concenpar}, the following corollary indicates that the above algorithm provides correct answers for approximate personalized \ced queries w.h.p.
\begin{corollary}
  Let $\theta = \Theta(\frac{\Lambda}{\epsilon k} \log{\frac{|V|}{p}})$ for any $\epsilon\in(0,1)$. Algorithm~\ref{alg:hindex} correctly answers an $\epsilon$-approximate personalized \ced query based on the h-index with probability at least $1 - p$.
\end{corollary}

\myp{Early Termination}
An early termination optimization can also accelerate h-index-based personalized \ced queries.
\begin{lemma}\label{lemma:hearlytermi}
  Given a query node $q$ and the quantile parameter $\lambda$, if there exists a node $u \in \vset^*$ such that $I(\bar{u}) > h(q)$ and $I(u) \geq \max(h(q), \lambda|V|)$, then the personalized \ced query returns \texttt{FALSE}.
\end{lemma}

\techreport{
\begin{proof}
  Each pair of nodes in $\iset(u)$ and $\iset(\bar{u})$ are neighbors of each other in $H$ according to the proof of Lemma~\ref{lemma:earlyterm}.
  Thus, each node $v_1 \in \iset(u)$ has at least $I(\bar{u})$ neighbors with degrees larger than or equal to $I(u)$.
  Since $I(\bar{u}) > h(q)$ and $I(u) \geq \max(h(q),\lambda\cdot|V|)$, there exist at least $I(u) \geq \lambda |V|$ nodes with h-indexes greater than or equal to $h(q)$.
  Therefore, $q$ is not a hub based on the h-index.
\end{proof}
}

Based on Lemma~\ref{lemma:hearlytermi}, \techreport{the sketch propagation can be terminated during \emph{Stage (I)} of the pivot-based partitioning method.
Specifically, we use the KMV sketches to estimate $I_f(u)$ and $I_b(u)$ for each $u\in\vset^*$.
If there exists a node $u \in \vset^*$ such that $\widetilde{I}_f(u) > N(q)$ and $\widetilde{I}_b(u)\geq\max(N(q), \lambda|V|)$, the algorithm can be terminated early since $N(q)$ is an upper bound of $h(q)$.}
once we obtain $\widetilde{h}(q)$ after \emph{Stage (I)}, the algorithm can be terminated without entering \emph{Stage (II)} if there exists a node $u \in \vset^*$ such that $\widetilde{I}(\bar{u}) > \widetilde{h}(q)$ and $\widetilde{I}(u) \geq \max(\widetilde{h}(q), \lambda|V|)$.

The following theorem bounds the probability that an early termination after \textit{Stage (I)} yields a false negative results for personalized \ced queries based on h-index.

\techreport{
\begin{theorem}\label{theorem:beta}
  Let $\theta = O(\frac{n}{\epsilon k}\log{\frac{|V|}{p}})$, the probability that a personalized \ced query based on h-index with answer \texttt{TRUE} is incorrectly answered with \texttt{FALSE} due to early termination during Stage (I) at node $u$ is at most $2\cdot p^{(\beta-1)/\epsilon}$, where $\beta = \min(\frac{\widetilde{I}_f(u)}{N(q)}, \frac{\widetilde{I}_b(u)}{\lambda|V|})$.
\end{theorem}

\begin{proof}
    According to Eqn.~\ref{equation:deltarelation}, for any $\delta<\beta-1$, we have 
    \begin{equation}
        I_f(u)\geq \frac{\widetilde{I}_f(u)}{1+\delta}\geq\frac{\beta N(q)}{1+\delta}>N(q)\geq h(q)
    \end{equation}
    and
    \begin{equation}
    \begin{aligned}
        I_b(u)&\geq\frac{\widetilde{I}_b(u)}{1+\delta}\geq\frac{\beta\cdot\max(N(q), \lambda|V|)}{1+\delta} \\ &>\max(N(q), \lambda|V|)\geq\max(h(q),\lambda|V|)
    \end{aligned}
    \end{equation}
    holds simultaneously with probability $1-2p^{(\beta-1)/\epsilon}$.

    Note that the early termination occurs at node $u$ leads to a false negative answer only if $\widetilde{I}_f(u) > h(q)$ and $\widetilde{I}_b(u) \geq \max(h(q),\lambda|V|)$, but $I_f(u) \leq h(q)$ or $I_b(u)< \max(h(q),\lambda|V|)$. Thus, a personalized \ced query with answer \texttt{TRUE} is incorrectly answered due to an early termination at $u$ with probability at most $2p^{(\beta-1)/\epsilon}$. 
\end{proof}
}

\begin{theorem}\label{theorem:beta:hindex}
  Let $\theta = O(\frac{\Lambda}{\epsilon k}\log{\frac{|V|}{p}})$, the probability that a personalized \ced query based on h-index with answer \texttt{TRUE} is incorrectly answered with \texttt{FALSE} due to early termination after Stage (I) at node $u$ is at most $(p+2p^{((1-\epsilon)\beta-1)/\epsilon})$, where $\beta = \min(\frac{\widetilde{I}(\bar{u})}{\widetilde{h}(q)}, \frac{\widetilde{I}(u)}{\lambda|V|})$.
\end{theorem}

\techreport{
\begin{proof}
    By Lemma~\ref{lemma:hq} and Eqn.~\ref{equation:BigOtransform}, for any $\delta<(1-\epsilon)\beta-1$,
    \begin{equation}
        I(\bar{u})\geq \frac{\widetilde{I}(\bar{u})}{1+\delta}\geq\frac{\beta \widetilde{h}(q)}{1+\delta}\geq \frac{\beta(1-\epsilon)h(q)}{1+\delta}>h(q)
    \end{equation}
    and
    \begin{equation}
    \begin{aligned}
        I(u)&\geq\frac{\widetilde{I}(u)}{1+\delta}\geq\frac{\beta\cdot\max(\widetilde{h}(q), \lambda|V|)}{1+\delta} \\ &\geq\frac{\beta\cdot\max((1-\epsilon)h(q), \lambda|V|)}{1+\delta} >\max(h(q), \lambda|V|)
    \end{aligned}
    \end{equation}
    holds with probability at least $1-2p^{((1-\epsilon)\beta-1)/\epsilon}-p$.

    Note that the early termination occurs at node $u$ leads to a false negative answer only if $\widetilde{I}(\bar{u}) \geq h(q)$ and $\widetilde{I}(u) > \max(h(q), \lambda|V|)$. However, we have $I(\bar{u}) < h(q)$ or $I(u)\leq \max(h(q), \lambda|V|)$. Thus, a personalized \ced query with answer \texttt{TRUE} is incorrectly answered due to an early termination at $u$ with probability at most $p+2p^{((1-\epsilon)\beta-1)/\epsilon}$. 
\end{proof}
}
